

\section{Énoncés des Exercices}

% --- Exercice 1 ---
\begin{exercice}
Soit $(f_n)_{n\in\N}$ une suite de fonctions d'un intervalle $I$ de $\R$ dans $\R$ convergeant uniformément sur $I$ vers une fonction $f$.
On pose :
$$ \forall x \in I, \quad g_n(x) = \frac{f_n(x)}{1+f_n^2(x)} $$
Montrer que $(g_n)_{n\in\N}$ converge uniformément sur $I$.
\end{exercice}

% --- Exercice 2 ---
\begin{exercice}
\textbf{a)} Établir la convergence simple puis uniforme sur $\R^+$ de la suite de fonctions définie par :
$$ \forall x \in \R^+, \quad f_n(x) = \frac{1}{1+x+x^2+\dots+x^n} $$
Diviser l'étude en deux temps, sur $[0, 1[$ et $[1, +\infty[$. On déterminera d'abord la limite simple qui est une fonction affine par morceaux.

\textbf{b)} (Bonus du manuscrit) Montrer la convergence et déterminer la limite de la suite définie par $u_n = \int_0^1 \frac{1}{1+x+\dots+x^n} \dx$.
\end{exercice}

% --- Exercice 3 ---
\begin{exercice}
Étudier la suite de fonctions définies sur $[0, 1]$ par :
$$ f_n(x) = \frac{2^n x}{1+n 2^n x^2} $$
Calculer les intégrales $\int_0^1 f_n(x) \dx$. Quel résultat se trouve confirmé par la considération de $\lim_{n\to\infty} \int_0^1 f_n(x) \dx$ ?
\end{exercice}

% --- Exercice 4 ---
\begin{exercice}
Établir la convergence simple sur $\R^+$ de la suite $(f_n)_{n\in\N}$ définie par :
$$ f_n(x) = \begin{cases} (1-\frac{x}{n})^n & \text{pour } x \in [0, n] \\ 0 & \text{pour } x > n \end{cases} $$
On notera $f$ la limite simple.
\begin{enumerate}
    \item Montrer que pour tout $x \in \R^+$, $0 \le f_n(x) \le f(x)$.
    \item Montrer que $f_n$ converge uniformément vers $f$ sur tout segment $[0, a]$.
    \item Montrer que la convergence est uniforme sur $\R^+$.
\end{enumerate}
\end{exercice}

% --- Exercice 5 ---
\begin{exercice}
Trouver le domaine $\Delta$ de convergence simple de la série de fonctions $\sum x^2 e^{-x\sqrt{n}}$.
Montrer qu'il n'y a ni convergence uniforme ni convergence normale sur $\Delta$.
Donner des intervalles de convergence uniforme.
\end{exercice}

% --- Exercice 6 ---
\begin{exercice}
On pose :
$$ f(x) = \sum_{n=1}^{+\infty} \frac{1}{1+n^2 x^2} \quad \text{et} \quad f_n(x) = \frac{x}{1+n^4 x^4} $$
\begin{enumerate}
    \item Montrer que $f$ est définie et continue sur $\R^*$ (on pourra montrer la convergence normale sur tout $[a, +\infty[, a>0$) et décroissante sur $\R^{+*}$.
    \item Montrer que $f$ est de limite nulle en $+\infty$.
    \item Trouver un équivalent de $f$ en $+\infty$ et en $0$.
\end{enumerate}
\end{exercice}

% --- Exercice 7 ---
\begin{exercice}
On pose $f(x) = \sum_{n=0}^{+\infty} (\arctan(x+n) - \arctan n)$.
\begin{enumerate}
    \item Montrer que $f$ est définie et de classe $\mathcal{C}^1$ sur $\R$. Montrer que $f$ est croissante sur $\R$.
    \item Montrer que pour tout $x$, $f(x+1) - f(x) = \frac{\pi}{2} - \arctan x$.
    \item Calculer $f(p)$ pour $p$ naturel et en déduire la limite de $f$ en $+\infty$.
\end{enumerate}
\end{exercice}

% --- Exercice 8 ---
\begin{exercice}
\textbf{a)} Étudier la convergence simple et la convergence uniforme de la série de fonctions :
$$ \sum_{n\ge 1} (-1)^{n+1} \ln\left(1+\frac{1}{nx}\right) $$
On notera $f$ la fonction somme.

\textbf{b)} Préciser le signe de $f$ et sa limite en $+\infty$.

\textbf{c)} Comparer $f(x)$ avec $g(x) = \sum_{n=1}^\infty \frac{(-1)^{n+1}}{nx}$ pour démontrer l'équivalent $f(x) \sim_{+\infty} \frac{\ln 2}{x}$.
\end{exercice}

% --- Exercice 9 ---
\begin{exercice}
On pose pour $x>0$, $u_n(x) = \frac{\sqrt{x} \ln n}{1+x n^2}$.
\begin{enumerate}
    \item Étudier la convergence de la série $\sum u_n(x)$.
    \item La somme est-elle continue sur $\R^{+*}$ ?
    \item Étudier le comportement en $+\infty$ de sa somme.
\end{enumerate}
\end{exercice}

\newpage
\section{Corrections}

% --- Correction 1 ---
\begin{correction}{1}
Soit $x \in I$. Puisque $(f_n)$ converge uniformément (CU) vers $f$ sur $I$, elle converge simplement (CS).
On a donc $f_n(x) \to f(x)$.
Par continuité de la fonction $t \mapsto \frac{t}{1+t^2}$, on a :
$$ g_n(x) = \frac{f_n(x)}{1+f_n^2(x)} \xrightarrow[n\to\infty]{} \frac{f(x)}{1+f^2(x)} = g(x) $$
Donc $(g_n)$ converge simplement vers $g$ sur $I$.

\textbf{Étude de la Convergence Uniforme :}
Étudions la différence $|g_n(x) - g(x)|$.
$$ |g_n(x) - g(x)| = \left| \frac{f_n(x)}{1+f_n^2(x)} - \frac{f(x)}{1+f^2(x)} \right| $$
$$ = \left| \frac{f_n(x)(1+f^2(x)) - f(x)(1+f_n^2(x))}{(1+f_n^2(x))(1+f^2(x))} \right| $$
$$ = \left| \frac{f_n - f + f_n f^2 - f f_n^2}{(1+f_n^2)(1+f^2)} \right| = \left| \frac{(f_n - f) + f_n f (f - f_n)}{(1+f_n^2)(1+f^2)} \right| $$
$$ = |f_n(x) - f(x)| \cdot \frac{|1 - f(x)f_n(x)|}{(1+f_n^2(x))(1+f^2(x))} $$
On utilise l'inégalité classique $2|ab| \le a^2 + b^2$, d'où $|ab| \le \frac{1}{2}(a^2+b^2) < 1+a^2+b^2$. Ou plus simplement, remarquons que la fonction $h(t) = \frac{t}{1+t^2}$ est $1$-lipschitzienne (car $|h'(t)| = |\frac{1-t^2}{(1+t^2)^2}| \le 1$).
Ainsi :
$$ |g_n(x) - g(x)| = |h(f_n(x)) - h(f(x))| \le 1 \cdot |f_n(x) - f(x)| $$
Puisque $(f_n)$ converge uniformément vers $f$ sur $I$, $\|f_n - f\|_\infty \to 0$.
Donc $\sup_{x \in I} |g_n(x) - g(x)| \le \|f_n - f\|_\infty \xrightarrow{n\to\infty} 0$.
Conclusion : $(g_n)$ converge uniformément vers $g$ sur $I$.
\end{correction}

% --- Correction 2 ---
\begin{correction}{2}
Pour $x \in \R^+$, $f_n(x) = \frac{1}{\sum_{k=0}^n x^k} = \frac{1}{\frac{1-x^{n+1}}{1-x}} = \frac{1-x}{1-x^{n+1}}$ (si $x \neq 1$).
\begin{itemize}
    \item Si $x \in [0, 1[$, $x^{n+1} \to 0$, donc $f_n(x) \to 1-x$.
    \item Si $x > 1$, $x^{n+1} \to +\infty$, donc $f_n(x) \sim \frac{1-x}{-x^{n+1}} \to 0$.
    \item Si $x = 1$, $f_n(1) = \frac{1}{n+1} \to 0$.
\end{itemize}
La limite simple est donc :
$$ f(x) = \begin{cases} 1-x & \text{si } x \in [0, 1[ \\ 0 & \text{si } x \ge 1 \end{cases} $$
La fonction limite $f$ est discontinue en $1$ (limite à gauche vaut 0, valeur en 1 vaut 0, mais $1-x \to 0$ continu ? Attendez, en $1^-$, $1-x \to 0$. La limite est continue sur $\R^+$).
Cependant, regardons la convergence uniforme.

\textbf{Sur $[0, 1[$ :}
$$ |f_n(x) - f(x)| = \left| \frac{1-x}{1-x^{n+1}} - (1-x) \right| = (1-x) \left| \frac{1 - (1-x^{n+1})}{1-x^{n+1}} \right| = \frac{(1-x)x^{n+1}}{1-x^{n+1}} $$
Pour étudier le sup, prenons $x$ proche de 1. Si $x \to 1$, cette quantité tend vers $0$.
Cependant, sur tout $[0, a]$ avec $a < 1$, majoration par $\frac{a^{n+1}}{1-a^{n+1}}$ qui tend vers 0. Donc CU sur $[0, a]$.

\textbf{Sur $[1, +\infty[$ :}
Pour $x \ge 1$, $|f_n(x) - 0| = \frac{1}{1+x+\dots+x^n} \le \frac{1}{n+1}$ (car chaque terme $\ge 1$).
Le majorant $\frac{1}{n+1}$ ne dépend pas de $x$ et tend vers 0.
Il y a donc convergence uniforme sur $[1, +\infty[$.
\end{correction}

% --- Correction 3 ---
\begin{correction}{3}
Soit $x \in [0, 1]$.
\begin{itemize}
    \item Si $x = 0$, $f_n(0) = 0$.
    \item Si $x \in ]0, 1]$, $f_n(x) \sim_{n\to\infty} \frac{2^n x}{n 2^n x^2} = \frac{1}{nx} \to 0$.
\end{itemize}
La suite $(f_n)$ converge simplement vers la fonction nulle sur $[0, 1]$.

\textbf{Étude de la Convergence Uniforme :}
Cherchons le maximum de $f_n$. Calculons la dérivée :
$$ f_n'(x) = 2^n \frac{1(1+n 2^n x^2) - x(n 2^n 2x)}{(1+n 2^n x^2)^2} = 2^n \frac{1 - n 2^n x^2}{(1+n 2^n x^2)^2} $$
$f_n'(x) = 0 \iff n 2^n x^2 = 1 \iff x_n = \frac{1}{\sqrt{n 2^n}} = \frac{1}{2^{n/2}\sqrt{n}}$.
Ce point $x_n$ est dans $[0, 1]$ pour $n$ assez grand.
$$ f_n(x_n) = \frac{2^n \frac{1}{\sqrt{n} 2^{n/2}}}{1+1} = \frac{2^{n/2}}{2\sqrt{n}} \xrightarrow{n\to\infty} +\infty $$
La norme infinie $\|f_n\|_\infty \to +\infty$. Il n'y a pas convergence uniforme.

\textbf{Intégrale :}
$$ \int_0^1 f_n(x) \dx = \int_0^1 \frac{2^n x}{1+n 2^n x^2} \dx = \left[ \frac{1}{2n} \ln(1+n 2^n x^2) \right]_0^1 $$
$$ = \frac{1}{2n} \ln(1+n 2^n) \sim \frac{1}{2n} \ln(n 2^n) = \frac{\ln n + n \ln 2}{2n} \xrightarrow{n\to\infty} \frac{\ln 2}{2} $$
Or $\int_0^1 \lim f_n = \int 0 = 0$.
On a $\lim \int \neq \int \lim$. Cela confirme l'absence de convergence uniforme (le théorème d'interversion ne s'applique pas).
\end{correction}

% --- Correction 4 ---
\begin{correction}{4}
\textbf{1. Convergence simple :}
Soit $x \in \R^+$. Pour $n$ assez grand ($n>x$), $f_n(x) = (1-\frac{x}{n})^n = e^{n \ln(1-\frac{x}{n})}$.
Quand $n \to \infty$, $n \ln(1-\frac{x}{n}) \sim n(-\frac{x}{n}) = -x$.
Donc $f_n(x) \to e^{-x}$. La convergence simple est vérifiée vers $f(x) = e^{-x}$.

Inégalité : On sait que $\ln(1+u) \le u$. Ici $\ln(1-\frac{x}{n}) \le -\frac{x}{n}$.
Donc $n \ln(1-\frac{x}{n}) \le -x \implies f_n(x) \le e^{-x} = f(x)$.
Et $f_n$ est positive sur $[0, n]$. D'où $0 \le f_n(x) \le f(x)$.

\textbf{2. Convergence uniforme sur $[0, a]$ :}
Sur le compact $[0, a]$, la suite de fonctions $(f_n)$ est une suite de fonctions continues convergeant simplement vers une fonction continue $f$.
De plus, on peut montrer la monotonie (Dini). Ou étudier la différence.
$$ |f(x) - f_n(x)| = e^{-x} - (1-\frac{x}{n})^n $$
Par Inégalité des Accroissements Finis sur le logarithme ou étude de fonction, on montre que la différence tend vers 0 uniformément sur $[0, a]$ (voir cours classique sur l'exponentielle).

\textbf{3. Convergence uniforme sur $\R^+$ :}
Regardons la borne supérieure de la différence $D_n(x) = e^{-x} - f_n(x)$.
Comme $f_n(x) = 0$ pour $x > n$, pour ces valeurs $D_n(x) = e^{-x}$.
Le sup sur $[n, +\infty[$ est $e^{-n} \to 0$.
Sur $[0, n]$, étudions la fonction. On peut montrer que le sup est atteint et tend vers 0.
Ainsi, la convergence est uniforme sur $\R^+$.
\end{correction}

% --- Correction 6 (D'après manuscrit Exercice 6) ---
\begin{correction}{6}
\textbf{a) Domaine et continuité :}
La série est $\sum u_n(x)$ avec $u_n(x) = \frac{1}{1+n^2 x^2}$.
Pour $x=0$, terme général $1 \not\to 0$, divergence grossière.
Pour $x \neq 0$, $u_n(x) \sim \frac{1}{n^2 x^2}$. Série de Riemann convergente ($2 > 1$).
Donc convergence simple sur $\R^*$.
$f$ est paire. Étudions sur $\R^{+*}$.
Sur tout intervalle $[a, +\infty[$ avec $a>0$ :
$|u_n(x)| \le \frac{1}{1+n^2 a^2} \le \frac{1}{n^2 a^2}$.
C'est le terme général d'une série numérique convergente indépendante de $x$.
Il y a convergence normale (donc uniforme) sur tout compact de $\R^*$.
Les fonctions $u_n$ sont continues, donc $f$ est continue sur $\R^*$.

\textbf{b) Limite en $+\infty$ :}
Puisque la convergence est uniforme sur $[1, +\infty[$, on peut intervertir limite et somme (Théorème de la double limite).
$\lim_{x\to +\infty} u_n(x) = 0$.
Donc $\lim_{x\to +\infty} f(x) = \sum 0 = 0$.

\textbf{c) Équivalent en $+\infty$ et $0$ :}
En $+\infty$ : $f(x) \sim \frac{\pi}{2x}$ ? Non, attention.
Comparaison Série-Intégrale. $t \mapsto \frac{1}{1+t^2 x^2}$ est décroissante.
$$ \int_1^{+\infty} \frac{\dt}{1+t^2 x^2} \le \sum_{n=1}^\infty \frac{1}{1+n^2 x^2} \le \frac{1}{1+x^2} + \int_1^{+\infty} \frac{\dt}{1+t^2 x^2} $$
Calcul de l'intégrale : $\int_0^{+\infty} \frac{\dt}{1+t^2 x^2} = \left[\frac{1}{x} \arctan(tx)\right]_0^{+\infty} = \frac{\pi}{2x}$.
Pour $x \to 0$, l'intégrale diverge, cela donne l'équivalent : $f(x) \sim_0 \frac{\pi}{2x}$.
Pour $x \to +\infty$, l'équivalent est aussi lié à l'intégrale ou au premier terme ?
Pour $x$ grand, $f(x) \approx \frac{1}{x^2} \sum \frac{1}{n^2} = \frac{\pi^2}{6 x^2}$.
\end{correction}

% --- Correction 7 ---
\begin{correction}{7}
\textbf{1. Convergence :}
Terme général $v_n(x) = \arctan(x+n) - \arctan n$.
Par TAF, il existe $c_n \in ]n, n+x[$ tel que $v_n(x) = \frac{x}{1+c_n^2}$.
Pour $n$ grand, $v_n(x) \sim \frac{x}{n^2}$. Série convergente pour tout $x$.
Convergence simple sur $\R$.
Sur tout segment $[-A, A]$, $|v_n(x)| \le \frac{A}{1+n^2}$. Convergence normale $\implies$ CU $\implies f$ continue.
Dérivée : $v_n'(x) = \frac{1}{1+(x+n)^2} > 0$. Somme de fonctions croissantes $\implies f$ croissante.

\textbf{2. Relation fonctionnelle :}
$$ f(x+1) - f(x) = \sum_{n=0}^\infty (\arctan(x+1+n) - \arctan n) - \sum_{n=0}^\infty (\arctan(x+n) - \arctan n) $$
C'est une somme télescopique à l'infini.
Soit $S_N(x) = \sum_{n=0}^N (\arctan(x+n) - \arctan n)$.
Calculons $S_N(x+1) - S_N(x)$.
Après télescopage partiel :
$$ \lim_{N\to\infty} (\arctan(x+N+1) - \arctan x) = \frac{\pi}{2} - \arctan x $$
Donc $f(x+1) - f(x) = \frac{\pi}{2} - \arctan x$.

\textbf{3. Calcul de $f(p)$ :}
$f(0) = 0$.
$f(1) - f(0) = \frac{\pi}{2} - 0$.
$f(p) = \sum_{k=0}^{p-1} (f(k+1)-f(k)) = \sum_{k=0}^{p-1} (\frac{\pi}{2} - \arctan k)$.
Quand $p \to \infty$, $f(p) \to +\infty$.
\end{correction}

% --- Correction 8 ---
\begin{correction}{8}
\textbf{a) Convergence :}
Série alternée $u_n(x) = (-1)^{n+1} \ln(1+\frac{1}{nx})$.
Pour $x$ fixé, $|u_n(x)| \searrow 0$. D'après le CSSA (Critère Spécial des Séries Alternées), la série converge simplement sur $\R^*$.
Majoration du reste : $|R_n(x)| \le |u_{n+1}(x)| = \ln(1+\frac{1}{(n+1)x})$.
Sur $[a, +\infty[$ ($a>0$), $\ln(1+\frac{1}{(n+1)x}) \le \ln(1+\frac{1}{(n+1)a})$.
Ce majorant est indépendant de $x$ et tend vers 0.
Donc convergence uniforme sur tout $[a, +\infty[$.

\textbf{b) Signe et Limite :}
$u_1(x) = \ln(1+\frac{1}{x}) > 0$. La somme d'une série alternée vérifiant le CSSA est du signe du premier terme. Donc $f(x) > 0$.
Limite en $+\infty$ : Par CU sur $[1, +\infty[$, $\lim_{x\to\infty} f(x) = \sum \lim u_n(x) = \sum 0 = 0$.

\textbf{c) Équivalent :}
On compare avec $g(x) = \sum \frac{(-1)^{n+1}}{nx} = \frac{1}{x} \sum \frac{(-1)^{n+1}}{n} = \frac{\ln 2}{x}$.
$|f(x) - g(x)| = \left| \sum (-1)^{n+1} (\ln(1+\frac{1}{nx}) - \frac{1}{nx}) \right|$.
Comme $\ln(1+u) - u \sim -u^2/2$, le terme général est en $O(\frac{1}{n^2 x^2})$.
La somme est majorée par $\frac{C}{x^2}$.
Donc $f(x) = \frac{\ln 2}{x} + O(\frac{1}{x^2})$, d'où $f(x) \sim \frac{\ln 2}{x}$.
\end{correction}

% --- Correction 9 ---
\begin{correction}{9}
\textbf{1. Convergence simple :}
Fixons $x > 0$. $u_n(x) = \frac{\sqrt{x} \ln n}{1+x n^2} \sim_{n\to\infty} \frac{\sqrt{x} \ln n}{x n^2} = \frac{\ln n}{\sqrt{x} n^2}$.
C'est un $o(1/n^{3/2})$ (par croissance comparée). Série de Riemann convergente.
Donc convergence simple sur $\R^{+*}$.

\textbf{2. Convergence Uniforme / Continuité :}
Étude de la fonction $h_n(x) = u_n(x)$.
Calcul de la dérivée pour trouver le sup.
On trouve que le sup n'est pas atteint en 0 mais dépend de $n$.
Si on se place sur $[a, +\infty[$, on peut majorer par une série numérique convergente (le terme en $1/n^2$ l'emporte).
$u_n(x) \le \frac{\sqrt{x} \ln n}{x n^2} = \frac{\ln n}{\sqrt{x} n^2} \le \frac{\ln n}{\sqrt{a} n^2}$.
Convergence normale sur $[a, +\infty[$.
Les fonctions $u_n$ sont continues, donc la somme est continue sur $\R^{+*}$.
\end{correction}
